\chapter{Introduction}
\label{chap:introduction}

\section{Background}

Remote sensing scene classification assigns semantic categories to aerial or
satellite image patches. It provides a means of organising image archives and
describing land use and land cover at the scene level. These descriptions can
support applications such as urban analysis and environmental monitoring,
although an image-level prediction does not locate every object or delineate
every land-cover boundary. Surveys by Zhu et al.~\cite{zhu2017deeprs} and
Cheng et al.~\cite{cheng2020scene} place scene classification within the wider
development of learned representations for remote sensing.

The target description depends on the task. A single-label benchmark assigns
one scene category to each image, even when several visual elements coexist.
NWPU-RESISC45 and the Aerial Image Dataset (AID) provide examples of this
setting~\cite{cheng2017nwpu,xia2017aid}. A multi-label task instead allows an
image to carry several land-cover labels, as in
BigEarthNet~\cite{sumbul2019bigearthnet,sumbul2021bigearthnetmm}. This
distinction changes the prediction objective and evaluation: choosing one
class differs from ranking and thresholding multiple label probabilities.
It also motivates examining regional representations in both settings,
without assuming that a useful representation for one will transfer
unchanged to the other.

Scene appearance presents a further challenge. Images within one category
may differ in scale, orientation, texture and composition, while different
categories can share visually similar components~\cite{cheng2017nwpu}.
Consequently, recognising local elements and combining their contextual
information are closely connected problems. An image representation must
retain useful evidence while reducing the many spatial features to the
compact vector required by a classifier.

Transfer learning offers a practical starting point. ImageNet established a
large natural-image resource~\cite{deng2009imagenet}, and Nogueira
et al.~\cite{nogueira2017exploiting} examined how pretrained convolutional
networks can be used for remote sensing scene classification. Domain-specific
pretraining provides another route: Seasonal Contrast uses temporal imagery
for self-supervised representation learning~\cite{manas2021seco}, while
SatMAE adapts masked pretraining to temporal and multispectral satellite
data~\cite{cong2022satmae}. These approaches illustrate that pretraining and
downstream architecture are separate design choices. This thesis holds the
recorded pretrained backbone source common within its internal comparison
and concentrates on the structure added after feature extraction.

\section{Motivation}

The shared backbone is Swin-Tiny (Swin-T), a hierarchical vision Transformer whose
shifted-window design enables interaction across neighbouring
windows~\cite{liu2021swin}. Its final tokens already contain contextual
information. The motivation for adding a graph is therefore to investigate
an additional regional inductive bias, rather than to remedy an assumed
absence of spatial relationships in the backbone. Whether that bias is
useful is an empirical question: an extra module can introduce helpful
structure, redundant computation or additional sensitivity to training.

Regional aggregation offers an intermediate representation between dense
feature tokens and a single global vector. Grouping tokens into a smaller
set of nodes makes their interactions explicit. Graph attention can then
aggregate neighbouring node features using learned
weights~\cite{velickovic2018gat}. However, this description leaves important
choices unresolved. Region membership, pooling weights, edge selection and
the return of graph features to the classifier each determine what
information is preserved and exchanged. An apparent improvement from a
complete architecture cannot automatically identify which choice caused it.

This thesis examines two grouping strategies within a Graph-Augmented Vision
Transformer (GAViT). The fixed strategy uses a deterministic spatial partition and
learns the relative importance of tokens within each region. The soft
strategy constructs image-dependent prototypes and permits contributions to
nearby anchors. Both retain a common anchor layout and explicit edge budget.
The comparison asks whether adaptive membership improves on a simple spatial
prior when the backbone and subsequent graph machinery are shared. It does
not equate attention-weighted pooling with learned segmentation, or assume
that soft assignment necessarily yields more meaningful regions.

The value of this structure also depends on its cost and consistency.
Parameter count, latency and memory describe different aspects of the added
computation. Repeated training can reveal sensitivity that a single favourable
score conceals. Visual inspection can show actual assignments, connections
and prediction changes, but cannot alone establish semantic correctness or
causal benefit. These considerations motivate evaluating the complete
accuracy--cost trade-off alongside diagnostic evidence.

\section{Research questions}

Three research questions organise the investigation:
\begin{enumerate}
 \item \textbf{RQ1: Regional graph reasoning.} Does fixed-partition GAViT
 improve classification relative to the shared Swin-T baseline, and how
 consistent is the observed difference across eligible training repetitions?
 \item \textbf{RQ2: Adaptive membership.} Does local prototype soft assignment
 improve on fixed grouping under the same task protocol, and what changes
 can be observed in regional aggregation, graph connectivity and predictions?
 \item \textbf{RQ3: Practical value and interpretation.} How do classification
 differences relate to additional computational cost, and what do successes,
 failures and visual diagnostics support about the role of the regional
 representation?
\end{enumerate}
These questions permit conditional and negative findings. A more flexible
grouping mechanism is not treated as preferable merely because it contains
additional learnable parameters.

\section{Contributions and scope}

The work makes three contributions. First, it implements a
Graph-Augmented Vision Transformer that combines attentive regional
aggregation, a sparse spatial-and-feature graph, graph attention and
region-to-token feedback. The architectural description distinguishes fixed
membership from learned pooling and distinguishes feature-based edge
selection from graph-attention weights. This makes the computational roles
of the components explicit.

Second, it formulates an image-dependent local prototype grouping mechanism
within the same architecture. Distinct normalisations define token membership
and regional pooling, while feedback uses the membership weights. This allows
the comparison to address adaptive aggregation without replacing the shared
backbone. The soft-versus-fixed
comparison concerns the resulting grouping pathway as a whole; it does not
separately identify every effect of the added projection and feedback map.

Third, it establishes an evaluation framework combining single-label and
multi-label classification, explicit model-selection rules, training
repetitions, computational measurements and qualitative diagnostics.
Calibration and final fitting are distinguished, and repeated training on
a frozen partition is distinguished from repeated random data splitting.
The contribution is the implemented design and its evidence-based
examination; consistent superiority over the baseline is a proposition to
be tested rather than an assumed outcome.

The empirical scope comprises red, green and blue (RGB) scene classification on NWPU-RESISC45,
AID and BigEarthNet-19. BigEarthNet is restricted to three optical bands in
this study, so the work does not evaluate full multispectral, radar or
multitemporal fusion. It also does not perform object detection or semantic
segmentation. Regional visualisations describe internal feature aggregation,
not independently validated object masks. Independent reconstruction of
HGTNet and DFAGCN forms a complementary comparison track; formal external
comparisons require the remaining implementation and protocol ambiguities
to be resolved.

\section{Organisation of the thesis}

The thesis is organised into seven chapters. Following this introduction,
Chapter~\ref{chap:related-work} reviews remote sensing Transformers, image graph networks and
regional grouping. Chapter~\ref{chap:method} presents GAViT and its fixed and
soft grouping variants. Chapter~\ref{chap:setup} specifies datasets,
preprocessing, model selection and evaluation. Chapter~5 reports
classification comparisons and component studies. Chapter~6 analyses
computational cost, qualitative evidence, failures and limitations.
Chapter~7 concludes the thesis and identifies directions for future work.
